\centerline{\bf \Huge Царь горы}


\begin{flushright}
        \scriptsize{Совершай ошибки. В такие моменты ты и учишься}
\end{flushright}


\problem{1}
Два кристалла и три слитка золота весят 8100 г, а три кристалла и три слитка золота — 9600 г. Сколько весит слитка золота?
% 1700 г

\problem{2}
Мистер Крабс задумал натуральное число, умножил на 3, стёр последнюю цифру, умножил на 2, вычел 4. Получилось число 10. Какое число мог задумать Мистер Крабс?
% 24, 25, 26

\problem{3}
Расставьте скобки и знаки арифметических действий так, чтобы получилось верное равенство: $4 \quad 5 \quad 9 \quad 8 \quad 7 \quad 10=2000$.
% 4⋅5⋅(9+8−7)⋅10=2000.

\problem{4}
Рико, Гавс и Спайк участвовали в легкоатлетическом забеге. В какой-то момент времени оказалось, что они (втроём) бегут рядом друг с другом, впереди них бежит половина участников забега и позади них — треть участников забега. Сколько спортсменов участвовало в забеге?
% 18


\problem{5}
Четверо ребят обсуждали ответ к задаче. 

Ваня сказал: «Это число 9».

Илья: «Это простое число».

Катя: «Это четное число».

А Маша сказала, что это число делится на 15.
Один мальчик и одна девочка ответили верно, а двое остальных ошиблись. Какой ответ в задаче на самом деле?
% 2

\problem{6}
В коробке шоколадные конфеты выложены в один слой в виде квадрата. Ранго съел все конфеты по периметру — всего 44 конфеты. Сколько конфет осталось в коробке?
% 100


\problem{7}
Сколькими способами число 2011 можно представить в виде суммы двух простых?
% 0

\problem{8}
10 часов тому назад от начала суток прошло столько же времени, сколько останется до конца суток через 2 часа. Сколько времени сейчас?
% 4 часа(дня/16:00)
